\ticket{1}{Парадигмы программирования}

\begin{examplan}
    \item Дать определение парадигмы программирования.
    \item Показать, что парадигма задает модель вычислений, структуру программы и способ декомпозиции задачи.
    \item Сравнить императивное, функциональное и логическое программирование.
    \item Для каждой парадигмы назвать теоретическую основу, ключевые понятия, достоинства, ограничения и примеры языков.
\end{examplan}

\subsection*{Короткая версия для письменного ответа}

\textbf{Парадигма программирования} --- это совокупность идей, понятий и методов, задающих фундаментальный стиль написания программ. Она определяет, как программист мыслит о задаче, как описывает вычисление и как организует программу.

Парадигма задает:
\begin{itemize}
    \item \textbf{модель вычислений}: последовательное изменение состояния, вычисление функций или логический вывод;
    \item \textbf{структуру программы}: процедуры, функции, классы, факты и правила;
    \item \textbf{основные строительные блоки}: переменные, объекты, функции, отношения;
    \item \textbf{способ декомпозиции задачи}: разбиение на шаги, преобразования данных или логические утверждения.
\end{itemize}

Парадигма не совпадает с языком программирования: многие современные языки являются мультипарадигменными. Например, Python поддерживает императивный, объектно-ориентированный и функциональный стиль.

\subsection*{Императивное программирование}

\textbf{Суть.} Вычисление описывается как последовательность команд, которые изменяют состояние программы. Программа отвечает на вопрос: \textit{как именно выполнить задачу}.

\textbf{Теоретическая и машинная основа.} Императивный стиль близок к машине Тьюринга и архитектуре фон Неймана: есть память, команды и пошаговое изменение содержимого памяти.

\textbf{Ключевые признаки:}
\begin{itemize}
    \item \textbf{состояние}: значения переменных и структур данных в конкретный момент времени;
    \item \textbf{присваивание}: операция, явно изменяющая состояние, например \texttt{x = x + 1};
    \item \textbf{управляющие конструкции}: последовательность команд, ветвления \texttt{if}, циклы \texttt{for}/\texttt{while};
    \item \textbf{побочные эффекты}: изменение внешнего состояния, например глобальных переменных, файлов, консоли или базы данных.
\end{itemize}

\textbf{Преимущества.} Императивный стиль интуитивен для алгоритмов, близок к реальной архитектуре компьютеров и часто позволяет эффективно управлять ресурсами.

\textbf{Недостатки.} В больших программах трудно контролировать изменяемое состояние. Побочные эффекты усложняют тестирование, доказательство корректности и параллельное выполнение.

\textbf{Примеры языков.} C, Pascal, Fortran; также C++, Java, Python и JavaScript поддерживают императивный стиль.

\subsection*{Функциональное программирование}

\textbf{Суть.} Вычисление рассматривается как вычисление значений выражений и применение функций. Программа ближе к описанию того, \textit{что нужно получить}, а не пошагового изменения памяти.

\textbf{Теоретическая основа.} Функциональное программирование опирается на лямбда-исчисление Алонзо Черча, где вычисление задается через абстракцию функций и применение функции к аргументу.

\textbf{Ключевые признаки:}
\begin{itemize}
    \item \textbf{чистые функции}: при одинаковых аргументах дают одинаковый результат и не имеют побочных эффектов;
    \item \textbf{ссылочная прозрачность}: вызов чистой функции можно заменить его результатом без изменения поведения программы;
    \item \textbf{неизменяемость данных}: вместо изменения существующего значения создается новое;
    \item \textbf{функции как объекты первого класса}: функции можно передавать как аргументы, возвращать из других функций и сохранять в переменных;
    \item \textbf{функции высшего порядка}: функции, работающие с другими функциями, например \texttt{map}, \texttt{filter}, \texttt{fold};
    \item \textbf{рекурсия}: часто используется вместо циклов;
    \item \textbf{ленивые вычисления}: в некоторых языках выражение вычисляется только тогда, когда его значение действительно нужно.
\end{itemize}

\textbf{Преимущества.} Чистые функции легче тестировать, комбинировать и распараллеливать. Неизменяемость данных уменьшает число ошибок, связанных с состоянием.

\textbf{Недостатки.} Для задач с большим количеством изменяемого состояния функциональный стиль может казаться менее естественным. Ввод-вывод и другие эффекты требуют специальных механизмов, например монад в Haskell. При неаккуратной реализации возможны накладные расходы на создание новых структур данных.

\textbf{Примеры языков.} Haskell, Lisp, Scheme, Clojure, OCaml, F\#; Scala и JavaScript являются мультипарадигменными языками с поддержкой функционального стиля.

\subsection*{Логическое программирование}

\textbf{Суть.} Программа задается как совокупность фактов и правил, а выполнение программы понимается как поиск доказательства запроса. Программист описывает, какие утверждения истинны и какие отношения существуют, а система вывода ищет значения переменных, при которых цель выполняется.

\textbf{Теоретическая основа.} Логическое программирование основано на логике предикатов первого порядка. На практике, например в Prolog, часто используется подмножество в виде дизъюнктов Хорна. Основные механизмы вывода --- унификация, резолюция и бэктрекинг.

\textbf{Ключевые признаки:}
\begin{itemize}
    \item \textbf{факты}: безусловные утверждения, например \texttt{parent(john, mary).};
    \item \textbf{правила}: условные утверждения, например \texttt{ancestor(X,Y) :- parent(X,Y).};
    \item \textbf{запросы}: цели, которые система пытается доказать;
    \item \textbf{отношения вместо функций}: одно отношение можно использовать в разных направлениях;
    \item \textbf{унификация}: подбор подстановок для переменных;
    \item \textbf{бэктрекинг}: возврат к альтернативам при неудаче или при поиске следующего решения.
\end{itemize}

\textbf{Пример идеи.} Если заданы факты \texttt{parent(john, mary).} и правило \texttt{ancestor(X,Y) :- parent(X,Y).}, то запрос \texttt{ancestor(john, Who)} ищет такое значение \texttt{Who}, при котором утверждение доказывается. В данном случае система найдет \texttt{Who = mary}.

\textbf{Преимущества.} Логический стиль дает высокий уровень абстракции и хорошо подходит для задач искусственного интеллекта, экспертных систем, символьных вычислений, обработки естественного языка и задач поиска.

\textbf{Недостатки.} Эффективность зависит от порядка правил, фактов и стратегии поиска. Для численных расчетов и низкоуровневого управления ресурсами логический стиль обычно менее удобен.

\textbf{Примеры языков.} Prolog, Mercury.

\subsection*{Сравнение парадигм}

\begin{description}
    \item[Императивная парадигма] описывает вычисление как изменение состояния командами. Основные элементы: переменные, присваивания, циклы, процедуры.
    \item[Функциональная парадигма] описывает вычисление как применение функций к значениям. Основные элементы: выражения, чистые функции, неизменяемые данные, функции высшего порядка.
    \item[Логическая парадигма] описывает вычисление как доказательство цели из фактов и правил. Основные элементы: предикаты, факты, правила, запросы, унификация.
\end{description}

Главное различие можно сформулировать так: императивная программа говорит компьютеру, \textit{как} получить результат; функциональная программа описывает, \textit{какое значение} нужно вычислить; логическая программа описывает, \textit{какие утверждения истинны}, а результат получается через логический вывод.

\begin{important}
    \item Парадигма --- это не синтаксис языка, а способ мышления о вычислении и структуре программы.
    \item Один язык может поддерживать несколько парадигм.
    \item Императивная парадигма: вычисление как последовательное изменение состояния.
    \item Функциональная парадигма: вычисление как применение функций и преобразование значений.
    \item Логическая парадигма: вычисление как доказательство запроса на основе фактов и правил.
    \item Для билета 1 достаточно объяснить унификацию, резолюцию и бэктрекинг на уровне идеи; подробный механизм Prolog относится к билету 4.
\end{important}

